\chapter{绪论}

\section{研究背景}

近年来，生成式人工智能（Artificial Intelligence Generated Content，AIGC）由文本生成逐步扩展到图像、音频、视频等多模态内容生产场景，正在改变数字内容的生产流程、表达方式与传播效率\cite{lit01}。在短视频平台高度普及的背景下，AI 配音、AI 绘图、AI 视频生成、虚拟人和场景重构等技术不断进入文旅内容生产过程，使地方文旅传播从传统实拍与人工剪辑，拓展到“真实素材+AI 加工”“AI 画面+文旅叙事”“AI 声音+目的地推荐”等多种形态。

文旅传播具有较强的视觉属性和体验属性。用户在真正到达目的地之前，往往先通过短视频、图文和平台评论形成对目的地的初步认知。既有研究表明，旅游短视频能够通过感知价值、情绪反应和目的地态度影响用户旅游意图\cite{lit04}，视觉叙事和内容吸引力也会进一步改变用户对目的地的想象方式\cite{lit06}。对于地方文旅部门而言，抖音等短视频平台已经不只是信息发布渠道，也是城市形象塑造、景区资源展示与用户互动的重要场域。

在这一背景下，AI 介入文旅短视频生产成为新的传播现象。一方面，AI 有助于降低内容生产门槛、提高视觉包装能力，帮助文旅账号更快生成具备想象力和可传播性的内容；另一方面，生成式内容也可能带来真实感不足、画面粗糙、风格同质化和用户信任下降等问题。OECD 关于人工智能与旅游的研究指出，AI 能够推动旅游创新与个性化体验，但也伴随着消费者信任、内容治理和目的地形象风险等问题\cite{lit23}。因此，AI 本身并不天然构成传播优势，其影响仍然需要结合具体内容特征加以检验。

从现实观察看，部分 AI 文旅短视频能够获得大量点赞、评论、收藏和转发，形成平台爆款；也有相当一部分视频虽然使用了明显的 AI 技术，但互动表现依然较弱。这说明，“是否用了 AI”并不足以解释传播效果差异。更值得研究的问题是：AI 的介入形式、介入程度、技术质感、视频风格和内容导向等因素，究竟哪些会显著影响传播效果？这些因素的作用方向和影响程度如何？基于此，本文以抖音平台省级文旅账号发布的 AI 文旅短视频为研究对象，围绕营销效果展开实证分析。

\section{文献综述与研究问题}

围绕本文主题，现有研究主要集中在 AIGC 内容生产、文旅短视频传播、社交媒体用户参与以及 AI 内容可信度等方向。AIGC 综述研究系统梳理了生成式人工智能在文本、图像、音频和视频等模态中的技术演进，指出 AIGC 正在推动内容生产从人工创作向人机协同创作转变\cite{lit01}。跨领域 AIGC 研究进一步强调，生成式内容已经进入数字营销、教育、公共健康等多类应用场景，并在内容生产效率和信息传递方式上产生显著影响\cite{lit26}。在视频内容评价方面，已有研究开始关注 AI 生成视频的视觉和谐性、文本一致性和整体质量，这为本文设置“AI 技术质感”变量提供了直接启发\cite{lit03}。

文旅短视频研究更多关注内容表达与用户旅游意图之间的关系。旅游短视频内容营销研究认为，短视频能够通过感知价值、情感反应和目的地态度影响用户旅游行为意图\cite{lit04}。关于 TikTok 短视频因素的研究发现，平台内容特征会影响不同代际用户的旅游行为意图\cite{lit05}。短视频视觉营销研究进一步表明，视觉内容和叙事吸引力会影响用户对目的地的感知和旅行意向\cite{lit06}。此外，也有研究从内容类型出发，说明旅游短视频能够通过激发潜在游客的想象体验影响出行意愿\cite{lit22}。这些成果说明，短视频已经成为旅游营销的重要媒介，但多数研究仍以一般旅游短视频为对象，对 AI 介入后的具体传播效果关注不足。

在传播效果测量方面，社交媒体用户参与研究为本文提供了重要参考。相关研究通常将点赞、评论、收藏、转发等行为视为用户对内容和品牌关系的外显反应，因此具备传播效果测量意义\cite{lit11}。与本文联系更紧密的是，AIGC 营销内容研究开始讨论生成内容如何通过价值感知和情感机制影响用户参与\cite{lit18}。这说明 AI 内容的传播效果并不只取决于技术是否被使用，还与用户对内容质量、价值与情绪体验的判断有关。

AI 内容真实感与可信度同样会影响用户接受。相关研究表明，用户会依据内容来源、表达质量和自身认知判断 AI 内容是否可信，感知可信度又会进一步影响态度与行为反应\cite{lit14}。在文旅场景中，这一问题更加突出。旅游内容不仅承担信息传播功能，也会影响用户对目的地的想象与信任。如果 AI 画面过度脱离地方真实特征，或者生成痕迹明显，就可能削弱内容的营销效果。

综合来看，已有文献分别讨论了 AIGC 内容生产、旅游短视频营销、社交媒体用户参与和 AI 内容可信度等问题，但仍存在进一步拓展空间。其一，已有旅游短视频研究较少将 AI 介入方式、介入程度和技术质感纳入同一分析框架。其二，现有 AIGC 研究偏重技术逻辑与应用设想，基于真实平台互动数据展开量化检验的成果仍较有限。其三，已有旅游传播研究更常使用旅游意图作为结果变量，而本文关注的是平台中可观察的点赞、评论、收藏和转发等互动行为。基于这些不足，本文尝试将 AIGC 特征、文旅内容特征和用户互动结果结合起来，分析 AI 文旅短视频营销效果的影响因素。

\subsection{研究问题}

本文围绕 AI 文旅短视频营销效果展开研究，核心问题是：在真实短视频平台中，哪些 AI 内容特征和视频内容特征会影响用户互动？本文从技术使用方式、画面呈现质量、视频主题、内容导向、视觉风格和发布时间条件等角度刻画样本特征，先比较不同类别之间的传播效果差异，再在控制多项因素后判断各指标是否仍具有显著影响。由此进一步回答：技术介入本身是否真正重要，还是画面质感和视频风格才是影响用户互动的关键。

\subsection{研究意义}

在理论意义上，本文将 AIGC 研究与文旅短视频传播研究结合起来，把技术使用方式、生成内容质量和视觉风格纳入旅游短视频传播效果分析框架，拓展了现有文旅短视频研究中对 AI 内容特征的关注。相较于多从旅游意图或内容吸引力切入的既有研究\cite{lit04}，本文进一步关注 AI 文旅内容在真实平台中的互动表现。

在实践意义上，本文能够为地方文旅部门和文旅账号优化 AI 短视频生产提供参考。当前不少账号容易把 AI 视为“新技术标签”，但真正能够转化为用户互动的，往往不是是否用了 AI，而是 AI 是否有效提升了内容质感、视觉记忆点和平台观看体验。研究结论可帮助运营者避免盲目追求 AI 介入深度，而将资源投入到画面质量、风格创新、视频节奏和用户感知体验等更关键的环节。

\subsection{研究创新点}

本文的创新点主要体现在以下三个方面。

第一，研究对象层面的创新在于，本文聚焦抖音平台省级文旅账号发布的 AI 文旅短视频。相较于一般旅游短视频或概念性 AIGC 研究，本文更关注官方文旅账号在真实平台环境中的 AI 内容传播实践，能够更直接反映地方文旅部门使用 AI 进行短视频营销的现实情况。

第二，变量体系层面的创新在于，本文没有简单使用“是否采用 AI”作为判断标准，而是进一步区分 AI 介入形式、AI 介入程度和 AI 技术质感，并结合视频主体、视频主题、内容导向和视频风格等变量展开分析。这样的设置有助于区分“用了 AI”和“AI 呈现得好不好”之间的差异。

第三，研究方法层面的创新在于，本文在内容分析和多元回归的基础上，引入稳健性检验、典型案例分析和预测性补充分析。回归分析用于识别主要影响因素，稳健性检验用于判断结论是否稳定，案例和机器学习方法则从具体内容和预测视角对结果进行补充说明。

\section{研究内容与方法}

本文的研究内容沿着“文献梳理、变量构建、差异识别、模型检验”的思路展开。研究先梳理 AIGC、文旅短视频、社交媒体用户参与和 AI 内容可信度相关文献，形成理论基础；再构建 AI 文旅短视频变量体系，对样本视频进行内容编码；随后使用描述统计和非参数检验分析不同变量类别之间的互动差异；在此基础上，通过多元回归、稳健性检验、典型案例补充讨论和预测性补充分析解释 AI 文旅短视频营销效果的影响因素。

本文主要采用内容分析法、描述统计、差异检验、多元 OLS 回归、稳健性检验和预测性补充分析。内容分析法用于根据视频内容和 AI 使用特征进行变量编码；描述统计和差异检验用于呈现样本分布并比较类别差异；多元 OLS 回归以 $\log(\text{总互动量}+1)$ 为因变量，同时纳入多个解释变量，识别各指标在控制其他因素后的净影响；高互动二分类 Logit 用于补充识别哪些内容特征更容易使视频进入高互动组；稳健性检验通过替换因变量、调整时间口径、处理极端值和加入省份固定效应判断结论稳定性；预测性补充分析则通过随机森林、梯度提升和高互动分类实验考察变量的预测价值，并作为解释性分析的补充。

本文将研究内容与方法概括为“理论基础构建--数据采集与筛选--变量编码--统计检验与回归分析--案例补充与机器学习验证--结论与建议”的总体框架。该框架依次呈现本文从理论基础、样本筛选、变量编码到实证检验、补充分析和结论建议的研究路径。具体而言，本文首先从 AIGC、文旅短视频传播、用户参与及 S-O-R 理论出发提出研究问题与分析框架；随后在样本层面完成候选视频汇总、AI 相关内容识别、样本筛选与变量编码；在实证层面依次展开描述统计、差异检验、多元回归与稳健性检验；最后结合典型案例分析和机器学习补充实验，对核心结论进行交叉印证，并据此形成研究结论与实践建议。图 \ref{fig:researchframework} 对这一研究内容与方法路径进行了概括。

\begin{figure}[!htbp]
  \centering
  \includegraphics[width=\textwidth]{../logs/extended_work/07_变量关系与实证分析路径图.png}
  \caption{研究内容与方法总体框架}
  \label{fig:researchframework}
\end{figure}

\section{论文结构}

本文共分为七章。第一章为绪论，介绍研究背景、文献综述、研究问题、研究意义、研究方法和论文结构。第二章为文献综述与理论基础，围绕 AIGC、文旅短视频、社交媒体用户参与、AI 真实感与信任以及 S-O-R 模型展开。第三章为研究设计，说明研究模型、研究假设、数据来源、变量定义和模型设定。第四章为实证分析，报告描述统计、差异检验、回归结果和稳健性检验。第五章为实验结果讨论与实践建议，围绕显著变量和非显著变量解释研究发现，并在结果讨论之后补充典型案例分析，进一步增强对定量结果的解释。第六章为预测性补充分析，报告不同机器学习模型的预测表现和变量重要性，用于检验变量体系的信息量和结论稳定性，而不是替代第三、四章的解释性模型。第七章为总结与展望，说明研究结论、研究不足和未来研究方向。
